Assume \(p\ge1\); zero arity is automatic. For rational profiles synchronize denominators and write \(u=-v(A)/q\), \(w=-v(B)/q\). Their maximum is \(-v(A+B)/q\). Take a common primary blowup making \(A,B,A+B\), the reference primary, and every other input ideal invertible, then normalize. Unit factors are harmless; the extra reference factor ensures an actual primary model even if other ideals are units. On this normal model put \(S=(A+B)\mathcal O_N\) and define ACTUAL ratio subideals
\[
L=A\mathcal O_N/S\subset\mathcal O_N,\qquad
M=B\mathcal O_N/S\subset\mathcal O_N,\qquad L+M=\mathcal O_N.
\]
They are invertible ideals, since all three original images are invertible. At an exceptional prime \(Z\) with \(u(\nu_Z)>w(\nu_Z)\), one has \(\operatorname{ord}_Z(A)<\operatorname{ord}_Z(B)\), so \(\operatorname{ord}_Z(M)>0\). Every local generator of \(M\) restricts to zero at the generic point of integral reduced \(Z\), hence to zero on its entire open restriction to \(Z\). Thus the entire map \(M|Z\to\mathcal O_Z\) is zero. Restricting \(L+M=\mathcal O_N\) gives a surjection \(L|Z\to\mathcal O_Z\), necessarily an isomorphism of line bundles. Therefore
\(\mathcal O_N(D_A-D_{A+B})|Z\) is trivial. In n4's independently geometric atomic formula, every mixed weight at an atom in the strict set is unchanged upon replacing the selected input \(u\) by \(\max(u,w)\). On the same common finite divisorial set, this proves
\[
\mu(\max(u,w),u_2,\ldots,u_p)|_{\{u>w\}}
=\mu(u,u_2,\ldots,u_p)|_{\{u>w\}}. \tag{MAX}
\]
This uses the actual subideals and their images, not just Weil coefficients or moment order.

For continuous \(C\) inputs approximate all entries uniformly by rational profiles. If a continuous test \(h\) has compact support \(K\) contained in the open set \(\{u>w\}\), the gap has a positive minimum on \(K\). Accurate approximants have the same strict inequality throughout \(K\). Apply (MAX) and n2's uniform weak continuity to pass the integrals against \(h\) to the limit. Thus (MAX) holds for \(C\) inputs. Now keep selected \(u,w\in C\) fixed and represent the other bounded \(D\) slots by actual decreasing \(C\)-nets. The strict open set and compactly supported test are fixed; n2's decreasing-net theorem proves (MAX) with all other slots in bounded \(D\).

Let bounded \(u,w\in D\) agree on an open \(U\), and fix the other bounded \(D\) slots. Choose actual decreasing \(C\)-representations \(u_\alpha\downarrow u\), \(w_\beta\downarrow w\), all with a common bound. Fix a continuous \(h\) with compact support \(K\subset U\), and \(\epsilon>0\). For each fixed \(\alpha\), the open sets
\[
\{w_\beta-\epsilon<u_\alpha\}
\]
increase and cover \(K\), because their pointwise limits there satisfy \(w-\epsilon=u-\epsilon<u_\alpha\). Compactness selects one \(\beta_0\). Its strict set is an open neighborhood \(V\) of \(K\), fixed for this \(\alpha,\epsilon\), and it remains strict for every \(\beta\ge\beta_0\). The continuous-selected-slot version of (MAX) gives
\[
\int h\,d\mu(\max(u_\alpha,w_\beta-\epsilon),\ldots)
=\int h\,d\mu(u_\alpha,\ldots).
\]
First decrease \(\beta\). The selected functions decrease to \(\max(u_\alpha,w-\epsilon)\), which belongs to bounded \(D\) by the actual representation or finite-max stability. N2's theorem applies with this prescribed limit. Then decrease \(\alpha\); the limit is the known bounded \(D\) function \(\max(u,w-\epsilon)\). Applying the same theorem again gives
\[
\int h\,d\mu(\max(u,w-\epsilon),\ldots)=\int h\,d\mu(u,\ldots). \tag{L1}
\]
These are two justified limits, not a flattening of nested nets. Interchange \(u,w\) to obtain the analogous identity for \(\max(w,u-\epsilon)\).

Both maximum functions are within \(\epsilon\) uniformly of the already known \(\max(u,w)\in D\). Although they increase as \(\epsilon\) decreases, n2's uniform estimate (U) applies INSIDE the established bounded domain. Letting \(\epsilon\) tend to zero in the two identities gives equality of the integrals for \(\mu(u,\ldots)\) and \(\mu(w,\ldots)\).

Such tests determine Radon restrictions to \(U\): compact inner regularity and compact-Hausdorff normality allow a continuous cutoff equal to one on any compact subset and supported in \(U\). Uniform approximation or the positive-functional description on these cutoffs then identifies the restricted measures. Hence \(\mu(u,\ldots)|U=\mu(w,\ldots)|U\). Replace slots successively to prove full mixed locality whenever corresponding tuples agree on \(U\). Every weak limit used a fixed continuous test with compact support, known domain membership and a common finite bound. No measure-order inference, increasing \(D\) closure, or plurifine claim occurs.
